\chapter{Introduction}
\label{ch:intro}

\section{Background}
\label{sec:intro-background}

Graph neural networks (GNNs) have become the dominant model family for collaborative filtering. Among them, LightGCN~\cite{he2020lightgcn} demonstrated that removing the non-linear activation and feature transformation from a GCN, and retaining only linear neighborhood aggregation, is sufficient to surpass deeper alternatives such as NGCF~\cite{wang2019ngcf} on standard recommendation benchmarks. The simplicity and stability of LightGCN have since established it as a strong, parameter-efficient backbone, on top of which more recent methods such as SGL~\cite{wu2021sgl} layer additional self-supervised objectives.

A complementary line of research, knowledge-graph-aware recommendation, augments collaborative signals with item-side semantics encoded in a knowledge graph (KG). Representative methods include KGAT~\cite{wang2019kgat}, which propagates messages over a unified user--item--entity graph; KGCL~\cite{yang2022kgcl} and MCCLK~\cite{zou2022mcclk}, which align KG-derived views with collaborative views via contrastive learning; and KGRec~\cite{yang2023kgrec}, which introduces edge-level rationalization to identify which KG triples are responsible for a recommendation. When the KG is dense and high-quality, these methods consistently outperform pure collaborative baselines.

\section{An Empirical Anomaly under Sparse KG}
\label{sec:intro-observation}

Our investigation begins with an empirical observation that motivates the rest of this thesis. On the review-derived KG dataset used throughout this work, an Amazon Books subset with on average $2.4$ aspect edges per item after filtering, all four representative KG-aware methods listed above are outperformed by pure LightGCN that does not consult the KG at all. Table~\ref{tab:intro-baseline} summarizes the NDCG@20 gap between these KG-aware methods and the no-KG baseline.

\begin{table}[H]
\centering
\caption{Comparison of KG-aware methods and pure LightGCN on NDCG@20.}
\label{tab:intro-baseline}
\begin{tabular}{lc}
\toprule
Method & NDCG@20 \\
\midrule
MCCLK~\cite{zou2022mcclk}      & 0.1067 \\
KGCL~\cite{yang2022kgcl}       & 0.1073 \\
KGAT~\cite{wang2019kgat}       & 0.1079 \\
KGRec~\cite{yang2023kgrec}     & 0.1095 \\
\midrule
Pure LightGCN~\cite{he2020lightgcn} (no KG) & \textbf{0.1179} \\
\bottomrule
\end{tabular}
\end{table}
\FloatBarrier

This is not a claim that the four baselines are flawed, each performs strongly on the dense KG datasets reported in their original papers. Rather, the observation suggests a failure mode shared by deep KG--CF coupling: when the KG is too sparse or noisy to provide reliable item semantics, fusing it directly into the scoring path drags KG noise into the gradient flow that would otherwise produce a clean collaborative signal.

\section{Sparse KG as a Real-World Default}
\label{sec:intro-sparse-kg}

A natural objection is to treat the empirical anomaly as a dataset peculiarity and either curate a denser KG, or apply KG completion to fill in missing edges. We argue that sparse KG is in fact the more common setting in practical recommendation, and that robustness under unreliable KG signal therefore deserves dedicated study rather than circumvention.

Three sources of sparsity recur in practice. First, \emph{review-derived KGs}, of the kind used in this work, inherit their density from whatever topics users happened to mention; coverage is inherently uneven. Second, \emph{cold-start and emerging domains} (new product categories, low-resource languages, niche verticals) rarely have access to a curated source comparable to Freebase or Wikidata, so the KG must be bootstrapped from limited signal. Third, \emph{privacy-constrained domains} such as medical or financial recommendation deliberately restrict which relational signals can be exposed to a model. In all three settings, an aggressive KG completion step is itself non-trivial: completion introduces its own noise, and typically requires sufficient seed signal to train, which is precisely what is missing in the sparse setting.

The implication is that a recommender's \emph{robustness} when the KG is unreliable is at least as important as its \emph{peak performance} when the KG is dense. This thesis takes that robustness as the primary design objective.

\section{Research Questions}
\label{sec:intro-rq}

The empirical anomaly and the practical importance of the sparse-KG setting motivate two research questions:

\begin{itemize}
  \item \textbf{(RQ1, diagnosis)} Why do existing KG-aware methods underperform pure LightGCN under sparse KG? What is the structural failure mode they share?
  \item \textbf{(RQ2, prescription)} What design principle allows a model to exploit the KG when its signal is informative, while preventing KG noise from contaminating collaborative filtering when it is not?
\end{itemize}

RQ1 is addressed in Chapters~\ref{ch:relatedwork} and~\ref{ch:experiments}, where we identify the shared failure mode as mandatory KG fusion under sparse and unreliable KG. RQ2 is developed in Chapter~\ref{ch:methodology} and validated empirically in Chapter~\ref{ch:experiments}, where we show that a gateable side-channel design can use KG signal when helpful and fall back to pure LightGCN when it is not.

The thesis statement of this work is that the KG should be treated as a gateable side channel rather than a mandatory component of the scoring pipeline. Chapter~\ref{ch:methodology} gives the concrete architectural and optimization choices that realize this principle.

\section{Thesis Organization}
\label{sec:intro-organization}

The remainder of this thesis is organized as follows. Chapter~\ref{ch:relatedwork} reviews related work along four axes: graph collaborative filtering, KG-aware recommendation, gating mechanisms in deep learning, and attention normalization. Chapter~\ref{ch:methodology} presents the RA-GARK architecture in full: it states the design principle, fixes notation and the problem formulation, and develops each module, the LightGCN local view, the global KG view with its two innovations (KG-SVD initialization and softmax aspect-saliency attention), the local-biased fusion gate, and the cross-view contrastive regularization, together with the training procedure and computational complexity. Chapter~\ref{ch:experiments} reports experimental results, ablations, sensitivity analyses, stability of results, and a case study. Chapter~\ref{ch:conclusion} discusses the methodological insight on attention normalization, acknowledges the limitations of single-dataset validation, and outlines future work.

\section{Contributions}
\label{sec:intro-contributions}

This thesis makes the following contributions:

\begin{enumerate}
  \item \textbf{An empirical study} demonstrating that, under sparse review-derived KG, four representative KG-aware methods (KGAT, KGCL, MCCLK, KGRec) are uniformly outperformed by pure LightGCN. We attribute this to the absence of an architectural mechanism that can disengage the KG when it is unreliable.

  \item \textbf{RA-GARK, a dual-view architecture} that integrates collaborative and KG-derived signals through a gateable side channel rather than deep fusion. Chapter~\ref{ch:methodology} gives the full architectural construction and training procedure.

  \item \textbf{Empirical validation} on the sparse review-derived KG benchmark, where RA-GARK surpasses both the strongest KG-aware baseline and pure LightGCN. The result shows that the proposed design principle flips the contribution of KG signal from \emph{net-negative} (as in the four baselines) to \emph{net-positive}, even when the KG itself is unchanged.

  \item \textbf{A methodological observation} regarding attention normalization in rationale-aware design. The effect of softmax versus sigmoid in our setting suggests that the choice of normalization should be aligned with the semantic assumption of the rationale. We discuss this as a single-dataset hypothesis rather than a general claim, pending replication on additional benchmarks (Chapter~\ref{ch:conclusion}).
\end{enumerate}
